\subsection{嘉当子代数的共轭性}

设 g 是一个李代数，h 是 g 的一个幂零子代数，R 是 h 在 g 中的非零权之集，换言之，是 h 上的线性形式 \(\lambda \neq 0\) 中使 \(\mathfrak{g}^{\lambda}(\mathfrak{h}) \neq 0\) 者所成之集，参看~§1 第 3 节命题 9 (iii)。假定

\[
\mathfrak {g} = \mathfrak {g} ^ {0} (\mathfrak {h}) \oplus \sum_ {\lambda \in R} \mathfrak {g} ^ {\lambda} (\mathfrak {h}),
\]

当 \(k\) 代数闭时即为此情形（§1，第 3 节，命题 9 (i)）。对 \(\lambda \in \mathbb{R}\) 和 \(x \in \mathfrak{g}^{\lambda}(\mathfrak{h})\) ，ad \(x\) 幂零（§1，第 3 节，命题 10 (iv)）。用 \(\mathrm{E}(\mathfrak{h})\) 记 \(\mathrm{Aut}_e(\mathfrak{g})\) 的由诸 \(e^{\mathrm{ad}x}\) （其中 \(x\) 形如上面）生成的子群。若 \(u \in \mathrm{Aut}(\mathfrak{g})\) ，立即有 \(u\mathrm{E}(\mathfrak{h})u^{-1} = \mathrm{E}(u(\mathfrak{h}))\) 。

引理 2. (i) 设 \(\mathfrak{h}_r\) 是 \(x\in \mathfrak{h}\) 中使 \(\mathfrak{g}^0 (x) = \mathfrak{g}^0 (\mathfrak{h})\) 者所成之集；这是 \(x\in \mathfrak{h}\) 中使 \(\lambda (x)\neq 0\) 对所有 \(\lambda \in \mathrm{R}\) 成立者所成之集，且 \(\mathfrak{h}_r\) 在 \(\mathfrak{h}\) 中按 Zariski 拓扑既开又稠密。

\begin{enumerate}
\def\labelenumi{(\roman{enumi})}
\setcounter{enumi}{1}
\tightlist
\item
  设 \(R = \{\lambda_{1}, \lambda_{2}, \ldots, \lambda_{p}\}\) ，其中诸 \(\lambda_{i}\) 两两不同。设 F 是从 \(\mathfrak{g}^{0}(\mathfrak{h}) \times \mathfrak{g}^{\lambda_{1}}(\mathfrak{h}) \times \cdots \times \mathfrak{g}^{\lambda_{p}}(\mathfrak{h})\) 到 g 的由公式
\end{enumerate}

\[
\mathrm{F} (h, x _ {1}, \dots , x _ {p}) = e ^ {\operatorname{ad} x _ {1}} \dots e ^ {\operatorname{ad} x _ {p}} h.
\]

定义的映射。那么 \(\mathrm{F}\) 是一个支配多项式映射（附录 I）。

断言 (i) 是明显的。我们证明 (ii)。设 \(n = \dim \mathfrak{g}\) 。若 \(\lambda \in \mathrm{R}\) 且 \(x \in \mathfrak{g}^{\lambda}(\mathfrak{h})\) ，我们有 \((\operatorname{ad} x)^n = 0\) 。由此推出 \((y, x) \mapsto e^{\operatorname{ad} x} y\) 是从 \(\mathfrak{g} \times \mathfrak{g}^{\lambda}(\mathfrak{h})\) 到 \(\mathfrak{g}\) 的多项式映射；用归纳法推出 \(\mathrm{F}\) 是多项式的。设 \(h_0 \in \mathfrak{h}_r\) ，并设 \(\mathrm{DF}\) 是 \(\mathrm{F}\) 在 \((h_0, 0, \ldots, 0)\) 处的切线性映射；我们证明 \(\mathrm{DF}\) 满射。对 \(h \in \mathfrak{g}^{0}(\mathfrak{h})\) ，有 \(F(h_{0} + h, 0, \ldots, 0) = h_{0} + h\) ，故 \(\mathrm{DF}(h, 0, \ldots, 0) = h\) ，从而 \(\mathrm{Im}(\mathrm{DF}) \supset \mathfrak{g}^{0}(\mathfrak{h})\) 。另一方面，对 \(x \in \mathfrak{g}^{\lambda_{1}}(\mathfrak{h})\) ，

\[
\mathrm{F} (h _ {0}, x, 0, \dots , 0) = e ^ {\operatorname{ad} x} h _ {0} = h _ {0} + (\operatorname{ad} x). h _ {0} + \frac {(\operatorname{ad} x) ^ {2}}{2 !} h _ {0} + \dots
\]

故 \(\mathrm{DF}(0,x,0,\ldots ,0) = (\mathrm{ad}x).h_0 = -(\mathrm{ad}h_0)x\) ；由于 \(\mathrm{ad}h_0\) 诱导 \(\mathfrak{g}^{\lambda_1}(\mathfrak{h})\) 的一个自同构，\(\operatorname {Im}(\mathrm{DF})\supset \mathfrak{g}^{\lambda_1}(\mathfrak{h})\) 。类似地，

\[
\operatorname{Im} (\mathrm{DF}) \supset \mathfrak {g} ^ {\lambda_ {i}} (\mathfrak {h})
\]

对所有 \(i\) 成立，故 DF 满射。于是附录 I 的命题 4 表明 F 是支配的。

命题 2. 假定 \(k\) 代数闭。设 \(\mathfrak{g}\) 是一个李代数，\(\mathfrak{h}\) 和 \(\mathfrak{h}'\) 是 \(\mathfrak{g}\) 的嘉当子代数。存在 \(u \in \mathrm{E}(\mathfrak{h})\) 和 \(u' \in \mathrm{E}(\mathfrak{h}')\) 使 \(u(\mathfrak{h}) = u'(\mathfrak{h}')\) 。

我们保留引理 2 的记号。由 \(\mathfrak{h}\) 和 \(\mathfrak{h}'\) 是嘉当子代数这一事实推出 \(\mathfrak{g}^0 (\mathfrak{h}) = \mathfrak{h}\) 且 \(\mathfrak{g}^0 (\mathfrak{h}') = \mathfrak{h}'\) 。由引理 2 和附录 I 的命题 3，\(\mathrm{E}(\mathfrak{h})\mathfrak{h}_r\) 和 \(\mathrm{E}(\mathfrak{h}')\mathfrak{h}_r'\) 包含 \(\mathfrak{g}\) 中按 Zariski 拓扑的稠密开子集。于是 \(\mathrm{E}(\mathfrak{h})\mathfrak{h}_r\cap \mathrm{E}(\mathfrak{h}')\mathfrak{h}_r' \neq \varnothing\) 。换言之，存在 \(u\in \mathrm{E}(\mathfrak{h}), u' \in \mathrm{E}(\mathfrak{h}')\) ，\(h\in \mathfrak{h}_r\) ，\(h' \in \mathfrak{h}_r'\) 使 \(u(h) = u'(h')\) ；那么

\[
u (\mathfrak {h}) = u \left(\mathfrak {g} ^ {0} (\mathfrak {h})\right) = \mathfrak {g} ^ {0} (u (h)) = \mathfrak {g} ^ {0} \left(u ^ {\prime} \left(h ^ {\prime}\right)\right) = u ^ {\prime} \left(\mathfrak {h} ^ {\prime}\right).
\]

推论. \(\mathrm{E}(\mathfrak{h}) = \mathrm{E}(\mathfrak{h}')\) 。

设 \(u, u'\) 如命题 2 中那样。那么

\[
\mathrm{E} (\mathfrak {h}) = u \mathrm{E} (\mathfrak {h}) u ^ {- 1} = \mathrm{E} (u (\mathfrak {h})) = \mathrm{E} (u ^ {\prime} (\mathfrak {h} ^ {\prime})) = u ^ {\prime} \mathrm{E} (\mathfrak {h} ^ {\prime}) u ^ {- 1} = \mathrm{E} (\mathfrak {h} ^ {\prime}),
\]

由此得到该推论。

由于这一结果，若 k 代数闭，我们就简单地用 E 表示群 \(\mathrm{E}(\mathfrak{h})\) ，其中 h 是 g 的嘉当子代数。

一般地，\(\operatorname{Aut}_{e}(\mathfrak{g}) \neq \operatorname{E}\) （例如，若 g 幂零，E 就化简为单位元，尽管只要 g 非交换就存在非平凡的初等自同构）。然而，可以证明（第八章，§10，习题 5），当 g 半单时 \(\operatorname{Aut}_{e}(\mathfrak{g}) = \operatorname{E}\) 。

定理 1. 假定 k 代数闭。设 g 是一个李代数。群 E 在 \(\operatorname{Aut}(\mathfrak{g})\) 中正规，且传递地作用在 g 的嘉当子代数所成之集上。

设 \(\mathfrak{h}\) 是 \(\mathfrak{g}\) 的嘉当子代数，\(v\in\mathrm{Aut}(\mathfrak{g})\) 。那么

\[
v \mathrm{E} (\mathfrak {h}) v ^ {- 1} = \mathrm{E} (v (\mathfrak {h})) = \mathrm{E} (\mathfrak {h}),
\]

故 \(\mathrm{E}(\mathfrak{h}) = \mathrm{E}\) 在 \(\mathrm{Aut}(\mathfrak{g})\) 中正规。若 \(h'\) 是 g 的另一个嘉当子代数，那么用命题 2 的记号，有 \(u'^{-1}u(\mathfrak{h}) = \mathfrak{h}'\) ，且 \(u'^{-1}u \in E\) 。

